\documentclass[11pt]{article}
\usepackage[margin=1in]{geometry}
\usepackage{amsmath,amssymb,amsthm,mathtools}
\usepackage{booktabs,longtable,array}
\usepackage{hyperref}
\usepackage{enumitem}
\usepackage{bm}
\usepackage{mathtools}
\usepackage[T1]{fontenc}
\usepackage{lmodern}
\setlength{\parskip}{0.5em}
\setlength{\parindent}{0pt}

\title{Extracted Notes on Clustered-Cyclic Codes, QUITS Usage,\
Seed Matrices, Decoding, and Logical Error Metrics}
\author{Prepared from the uploaded manuscript}
\date{\today}

\begin{document}
\maketitle

\section{Purpose}
This note extracts, reorganizes, and condenses from the uploaded manuscript the following items:
\begin{enumerate}[label=(\arabic*)]
    \item the clustered-cyclic (CC) code construction,
    \item how \texttt{QUITS} is used in the circuit-level memory simulations and what is done differently,
    \item the seed matrices for the $[24,8,3]$, $[40,8,5]$, and $[56,8,7]$ CC codes,
    \item the decoding settings used in the simulations,
    \item how the logical error / failure quantities are computed from repeated-syndrome simulations.
\end{enumerate}
Throughout, the notation follows the manuscript as closely as possible.

\section{Clustered-cyclic code construction}

\subsection{Base ring and seed matrices}
A clustered-cyclic code is defined as a lifted-product (LP) CSS code over the quotient ring
\[
R = \mathbb{F}_2[x]/(x^p+1),
\]
where the lift parameter $p$ is prime. The construction starts from two \emph{square} classical seed matrices
\[
H_a \in R^{n_a\times n_a}, \qquad H_b \in R^{n_b\times n_b}.
\]
The manuscript imposes the following structural conditions on these seed matrices:
\begin{itemize}
    \item every entry is either a \emph{binomial} or zero,
    \item the supports of the nonzero monomials are arranged cyclically,
    \item every row and every column has the same number of nonzero binomials,
    \item both seed matrices are full rank.
\end{itemize}

\subsection{CSS parity-check matrices}
Given $H_a$ and $H_b$, the CC code parity-check matrices are
\begin{align}
H_X &= \left(\begin{array}{c|c} I_{n_b} \otimes H_a & H_b^{*} \otimes I_{n_a} \end{array}\right), \\
H_Z &= \left(\begin{array}{c|c} H_b \otimes I_{n_a} & I_{n_b} \otimes H_a^{*} \end{array}\right),
\end{align}
where $H_a^{*}$ and $H_b^{*}$ denote conjugate-transposes over the ring, and $I_m$ is the $m\times m$ identity over $R$.

\subsection{Parameters}
The manuscript states that the resulting binary CSS code has parameters
\[
N = 2pn_an_b,
\qquad
k = 2n_an_b,
\qquad
d \le p.
\]
If the row weights of $H_a$ and $H_b$ are $w_a$ and $w_b$, then the stabilizer check weight is
\[
W = 2(w_a+w_b).
\]
For the examples extracted below, the paper focuses on two families:
\begin{itemize}
    \item $n_a=n_b=2$, giving $k=8$ logical qubits,
    \item $n_a=n_b=3$, giving $k=18$ logical qubits.
\end{itemize}

\subsection{Why the code is called ``clustered-cyclic''}
Let
\[
\chi = 1+x+x^2+\cdots+x^{p-1} \in R.
\]
A key theorem in the manuscript shows that CC codes admit a \emph{clustered logical operator basis}. In the ring description, the logical basis matrix is
\[
L=
\left(
\begin{array}{c|c}
\operatorname{diag}_{n_an_b}(\chi) & 0 \\
\hline
0 & \operatorname{diag}_{n_an_b}(\chi)
\end{array}
\right).
\]
After passing to the binary representation, each occurrence of $\chi$ becomes a support on a block of $p$ physical qubits. Hence each logical operator is supported on one entire \emph{cluster} of size $p$, and each logical $X$ operator anti-commutes with exactly one logical $Z$ operator supported on the same cluster. This is the structural feature exploited later for parallel logical measurements and surgery scheduling.

\subsection{What is special relative to more generic LP / HGP constructions}
The manuscript's extracted differences are:
\begin{enumerate}[label=(\alph*)]
    \item the seed matrices are restricted to a cyclic binomial/zero pattern rather than being arbitrary LP seeds,
    \item the seeds are square and full-rank,
    \item the resulting logical basis is explicitly clustered, so logical representatives are highly structured and directly addressable,
    \item this clustered basis is what enables the parallel product-surgery constructions used later in the paper.
\end{enumerate}
In other words, the paper is not simply using an arbitrary lifted-product code; it is carving out a subfamily whose logical basis is engineered to make logical measurement layers more surface-code-like.

\section{How \texttt{QUITS} is used, and what is done differently}

\subsection{Role of \texttt{QUITS}}
For circuit-level memory simulations, the manuscript states that the repeated-syndrome simulation framework is \emph{built upon and extends} \texttt{QUITS}. In each syndrome extraction round, data qubits are coupled to auxiliary qubits by CNOT gates determined by the Tanner graph of the code.

Thus, at a high level, \texttt{QUITS} is used as the base infrastructure for repeated-syndrome circuit-level simulation of qLDPC memory experiments.

\subsection{What is modified relative to the generic \texttt{QUITS} scheduling}
The manuscript describes a direction-aware greedy edge-coloring scheduler on the Tanner graph. For this task, the experimental conditions in \texttt{experiment.md} fix the physical schedule using DSATUR on QUITS v1.0.0's \texttt{cardinal} layout. The motivation is to reduce the CNOT depth per syndrome round, because every additional layer also adds idle time for qubits that are not active in that layer.

The extracted custom scheduling procedure is:
\begin{enumerate}[label=(\arabic*)]
    \item \textbf{Directional partitioning.} Partition Tanner-graph edges into four disjoint groups
    \[
    v\in\{N,S,E,W\},
    \]
    interpreted as distinct interaction directions. The manuscript notes that this is natural for HGP/LP-style constructions where two-qubit gates can be viewed as horizontal or vertical Tanner-graph edges.

    \item \textbf{Conflict graph construction.} For each direction $v$, form a conflict graph $\mathcal G_v$ whose vertices are CNOT gates in that direction, and connect two vertices exactly when the corresponding CNOTs share a qubit.

    \item \textbf{DSATUR graph coloring.} Use NetworkX DSATUR with the vertex ordering stated in \texttt{experiment.md}. Each color class gives a set of pairwise non-conflicting CNOTs that can be executed simultaneously.

    \item \textbf{Circuit assembly.} The CNOT depth is taken to be
    \[
    \gamma_{\mathrm{CNOT}} = \sum_v \phi(\mathcal G_v),
    \]
    where $\phi(\mathcal G_v)$ denotes the chromatic number returned by the coloring procedure. In addition, the Hadamards around the $X$-type ancillas and the mid-round ancilla measurement/reset stages each contribute one extra time-step.
\end{enumerate}

\subsection{Bottom line}
So the paper's use of \texttt{QUITS} is:
\begin{itemize}
    \item use \texttt{QUITS} as the base repeated-syndrome simulation framework,
    \item use the direction-aware DSATUR Tanner-graph edge-coloring schedule defined in the experimental conditions,
    \item which removes the circuit-order ambiguity that otherwise changes absolute logical-failure rates.
\end{itemize}
The manuscript does \emph{not} give further implementation details beyond that high-level description, so this is the cleanest faithful extraction.

\section{Decoding and post-processing settings}
The paper states that syndrome-extraction circuits are compiled and sampled with \texttt{Stim}. Decoding is then performed by BP+OSD on the \emph{full space-time parity-check matrices} derived from the detector error model (DEM).

For this task, \texttt{experiment.md} defines the decoder's fault-variable
representation. Each flattened DEM error instruction remains a separate
variable with its full detector and logical-observable effects and its prior;
detector-set-only aggregation is not used.

The reported hyperparameters are:

\begin{center}
\begin{tabular}{lc}
\toprule
Parameter & Value \\
\midrule
BP update rule & Min-sum \\
Min-sum scaling factor & 0.625 \\
Maximum BP iterations & 200 \\
OSD method & OSD-CS \\
OSD order & 5 \\
\bottomrule
\end{tabular}
\end{center}

So, in compact form, the decoding stack is:
\[
\texttt{Stim} \;\longrightarrow\; \text{DEM} \;\longrightarrow\; \text{full space-time parity-check matrix} \;\longrightarrow\; \text{BP+OSD}.
\]

\section{How the logical error / failure quantities are computed}
The manuscript's appendix for circuit-level memory simulations works with repeated syndrome extraction over $\eta$ rounds.

\subsection{Choice of the number of rounds}
The number of rounds is chosen as
\[
\eta = 42 \quad \text{if } \epsilon < 4\times 10^{-3},
\qquad
\eta = 20 \quad \text{if } \epsilon > 4\times 10^{-3},
\]
where $\epsilon$ is the physical error rate.

\subsection{From total logical failure over $\eta$ rounds to per-round LFR}
Let $\epsilon_\gamma$ denote the empirically estimated \emph{total logical failure rate over $\eta$ rounds}. Assuming independent failures across rounds in the low-error regime, the per-round logical failure rate (LFR) is defined through
\[
(1-\mathrm{LFR})^{\eta} = 1-\epsilon_\gamma.
\]
Therefore,
\[
\boxed{\mathrm{LFR} = 1-(1-\epsilon_\gamma)^{1/\eta}}.
\]

\subsection{Per-logical-qubit, per-round quantity}
To compare codes with different numbers of logical qubits $k$, the manuscript further defines an effective per-logical-qubit per-round failure rate $\epsilon_{1,1}$ by assuming independent, uniform failures across the $k$ logical qubits:
\[
1-\epsilon_\gamma = (1-\epsilon_{1,1})^{k\eta}.
\]
Hence,
\[
\boxed{\epsilon_{1,1} = 1-(1-\epsilon_\gamma)^{1/(k\eta)}}.
\]

\subsection{Interpretation}
So there are really three related quantities in play:
\begin{enumerate}[label=(\roman*)]
    \item $\epsilon_\gamma$: total logical failure probability over the full $\eta$-round experiment,
    \item LFR: logical failure rate per syndrome cycle,
    \item $\epsilon_{1,1}$: effective per-logical-qubit, per-cycle failure rate.
\end{enumerate}
If by ``LER'' one means the logical error/failure rate reported in the circuit-level appendix, the paper computes it through this pipeline.

\section{Extracted seed matrices}
In all three examples below, the seed matrices are given over
\[
R=\mathbb{F}_2[x]/(x^p+1),
\]
with the indicated prime lift parameter $p$.

\subsection{$[24,8,3]$ code, with $p=3$}
This is a $k=8$ CC code, so $n_a=n_b=2$ and the seed matrices are $2\times2$:
\[
H_a=
\begin{pmatrix}
1+x^2 & x+x^2 \\
1+x   & x+x^2
\end{pmatrix},
\qquad
H_b=
\begin{pmatrix}
1+x^2 & x+x^2 \\
1+x   & 1+x
\end{pmatrix}.
\]

\subsection{$[40,8,5]$ code, with $p=5$}
Again $n_a=n_b=2$:
\[
H_a=
\begin{pmatrix}
x+x^3 & x+x^4 \\
x+x^4 & x+x^4
\end{pmatrix},
\qquad
H_b=
\begin{pmatrix}
1+x   & x^3+x^4 \\
1+x^4 & 1+x^4
\end{pmatrix}.
\]

\subsection{$[56,8,7]$ code, with $p=7$}
This is a $k=8$ CC code, so $n_a=n_b=2$ and the seed matrices are $2\times2$:
\[
H_a=
\begin{pmatrix}
x+x^5 & x+x^3 \\
1+x^2 & 1+x^3
\end{pmatrix},
\qquad
H_b=
\begin{pmatrix}
x^2+x^5 & 1+x \\
x^5+x^6 & x^4+x^6
\end{pmatrix}.
\]

\section{Compact summary}
A faithful extraction from the manuscript is the following.

\begin{itemize}
    \item A clustered-cyclic code is a special lifted-product CSS code over $R=\mathbb{F}_2[x]/(x^p+1)$ built from square full-rank cyclic binomial seed matrices $H_a,H_b$.
    \item Its stabilizer matrices are
    \[
    H_X = (I_{n_b}\otimes H_a\mid H_b^*\otimes I_{n_a}),
    \qquad
    H_Z = (H_b\otimes I_{n_a}\mid I_{n_b}\otimes H_a^*).
    \]
    \item The code has parameters $N=2pn_an_b$, $k=2n_an_b$, and $d\le p$, and admits a clustered logical basis supported on blocks of $p$ qubits.
    \item Circuit-level memory simulations use \texttt{QUITS}; the task fixes the directional edge-coloring implementation to the DSATUR schedule defined in the experimental conditions.
    \item Syndrome-extraction circuits are compiled and sampled with \texttt{Stim}; decoding is BP+OSD on the full space-time parity-check matrix derived from the DEM.
    \item The simulation post-processing turns the measured total logical failure probability over $\eta$ rounds into a per-round LFR and then into a per-logical-qubit, per-round quantity $\epsilon_{1,1}$.
\end{itemize}

\end{document}
